\documentclass[a4paper,11pt]{article}

\usepackage[margin=2cm]{geometry}
\usepackage{graphicx}
\usepackage{float}
\usepackage{booktabs}
\usepackage{longtable}
\usepackage{array}
\usepackage{xcolor}
\usepackage{hyperref}
\usepackage{enumitem}

\setcounter{secnumdepth}{0}

\hypersetup{
	colorlinks=true,
	linkcolor=blue,
	urlcolor=blue
}

\title{Decision on Type Validation and Typed Loading}
\author{Giuseppe Rudi}
\date{\today}

\begin{document}
	
	\maketitle

	
	During the design of the Data Quality Assessment and Data Cleaning workflow, two possible strategies were considered.
	
	The first strategy was to load all re-engineered CSV files into a staging PostgreSQL schema where every column was stored as text.
	This solution would allow the DQA to explicitly test whether each value can be converted into the expected type, such as date, boolean, integer, or floating-point number.
	
	However, this strategy is more suitable for uncontrolled CSV sources, where all values are initially textual and the project must infer and validate the correct data type of each column.
	
	This is not the main situation of the present project.
	
	In this project, the main data source is the FastF1 Python API.
	FastF1 does not provide uncontrolled textual CSV files as the primary input.
	Instead, it returns structured Pandas DataFrames, where columns already have an expected computational type or a structured representation.
	
	During the re-engineering phase, the project selects the relevant attributes, adds structural identifiers, converts time-related attributes into integer milliseconds, and exports coherent reconciled CSV files.
	
	For this reason, the project adopts a typed loading strategy.
	
	The re-engineered CSV files are loaded into PostgreSQL tables with explicit data types.
	This choice is not used as a replacement for data quality assessment.
	Rather, it reflects the fact that basic type interpretation has already been handled upstream by the API and by the controlled Python re-engineering pipeline.
	
	Therefore, the post-load DQA does not focus on pure type parseability checks such as:
	
	\begin{itemize}
		\item whether a timestamp column can be parsed as a timestamp;
		\item whether a boolean column can be parsed as a boolean;
		\item whether a numerical column can be parsed as a number.
	\end{itemize}
	
	These checks would be more relevant if the project started from uncontrolled textual files.
	In this project, they would add complexity without providing significant analytical value.
	
	Instead, the DQA focuses on the quality problems that are still possible and relevant after typed loading.
	
	A value may have the correct technical type but still be invalid, inconsistent, suspicious, or analytically unusable.
	
	For example:
	
	\begin{itemize}
		\item a speed value can be numerical but negative or unrealistic;
		\item a lap time can be numerical but equal to zero or inconsistent with sector times;
		\item two dates can be valid timestamps but logically inconsistent with each other;
		\item a tyre compound can be textual but outside the expected Formula 1 compound domain;
		\item a row can have valid types but missing identifiers;
		\item two rows can be individually valid but duplicated according to a natural key;
		\item a foreign key can refer to a non-existing parent row;
		\item a lap can be valid as data but not suitable for clean performance analysis.
	\end{itemize}
	
	For this reason, the project distinguishes between two levels of validity:
	
	\begin{itemize}
		\item \textbf{type-level validity}, which concerns whether a value can be represented with the expected technical type;
		\item \textbf{domain-level and semantic validity}, which concerns whether a value is meaningful, coherent, and usable in the Formula 1 analytical domain.
	\end{itemize}
	
	The first level is mostly handled by the source API, the re-engineering script, and the typed loading process.
	The second level is the main focus of the Data Quality Assessment and Data Cleaning phases.
	
	Examples of post-load DQA checks are:
	
	\begin{itemize}
		\item missing required identifiers;
		\item duplicated primary or natural keys;
		\item invalid categorical domains;
		\item impossible or unrealistic numerical values;
		\item inconsistencies between related attributes;
		\item mismatches between related tables;
		\item suspicious or unreliable laps;
		\item referential integrity problems.
	\end{itemize}
	
	This decision keeps the DQA and Data Cleaning phases focused on the problems that are most relevant for the construction of the Data Warehouse.
	
	In particular, the project does not attempt to duplicate basic type inference already provided by the API and by the re-engineering pipeline.
	Instead, it evaluates whether the typed and loaded values are complete, consistent, realistic, referentially correct, and analytically useful.
	
\end{document}